% ============================================================
\chapter*{Pr\'eface}
\addcontentsline{toc}{chapter}{Pr\'eface}
% ============================================================

La \textbf{recherche op\'erationnelle} (RO) est une discipline scientifique qui
s'appuie sur des m\'ethodes math\'ematiques, statistiques et algorithmiques pour
aider \`a la prise de d\'ecision. N\'ee pendant la Seconde Guerre mondiale au sein
des \'etats-majors alli\'es, elle s'est depuis impos\'ee comme un outil
incontournable dans l'industrie, la logistique, la finance, les
t\'el\'ecommunications et bien d'autres domaines.

\section*{Objectifs du cours}

Ce cours s'adresse aux \'etudiants de deuxi\`eme et troisi\`eme ann\'ee de licence
(L2--L3) en math\'ematiques, informatique ou sciences de l'ing\'enieur.
\`A l'issue de cet enseignement, l'\'etudiant sera capable de~:

\begin{itemize}
  \item Mod\'eliser des probl\`emes concrets sous forme de programmes
        math\'ematiques (lin\'eaires, en nombres entiers, non lin\'eaires).
  \item R\'esoudre ces programmes \`a l'aide des algorithmes fondamentaux
        (simplexe, branch-and-bound, plus court chemin, flot maximal, etc.).
  \item Interpr\'eter les r\'esultats, notamment via l'analyse de dualit\'e et de
        sensibilit\'e.
  \item Utiliser des outils logiciels (Python, PuLP, SciPy, NetworkX) pour
        r\'esoudre des probl\`emes de taille r\'ealiste.
  \item Comprendre les principes de la simulation \`a \'ev\'enements discrets et de
        la th\'eorie des files d'attente.
\end{itemize}

\section*{Organisation de l'ouvrage}

L'ouvrage est organis\'e en onze chapitres, con\c{c}us pour \^etre \'etudi\'es de
mani\`ere s\'equentielle~:

\begin{description}
  \item[Chapitre 1 -- Introduction \`a la RO.]
    Historique, d\'efinitions et m\'ethodologie g\'en\'erale de la mod\'elisation.
  \item[Chapitre 2 -- Programmation lin\'eaire~: formulation et g\'eom\'etrie.]
    Construction de mod\`eles lin\'eaires, interpr\'etation g\'eom\'etrique, r\'egions
    admissibles et sommets optimaux.
  \item[Chapitre 3 -- M\'ethode du simplexe.]
    L'algorithme du simplexe sous forme tabulaire, initialisation par la
    m\'ethode du grand $M$ et la m\'ethode \`a deux phases.
  \item[Chapitre 4 -- Dualit\'e en programmation lin\'eaire.]
    Construction du dual, th\'eor\`emes de dualit\'e, interpr\'etation \'economique.
  \item[Chapitre 5 -- Analyse de sensibilit\'e.]
    Variation des coefficients, analyse param\'etrique, interpr\'etation des
    prix fictifs (\emph{shadow prices}).
  \item[Chapitre 6 -- Probl\`emes de transport et d'affectation.]
    Mod\`eles de transport \'equilibr\'es et non \'equilibr\'es, m\'ethode du
    stepping-stone, probl\`eme d'affectation et algorithme hongrois.
  \item[Chapitre 7 -- Th\'eorie des graphes appliqu\'ee.]
    Plus courts chemins (Dijkstra, Bellman-Ford), arbres couvrants minimaux
    (Kruskal, Prim).
  \item[Chapitre 8 -- Flots dans les r\'eseaux.]
    Flot maximal (Ford-Fulkerson), coupe minimale, flot \`a co\^ut minimal.
  \item[Chapitre 9 -- Programmation en nombres entiers.]
    M\'ethode de branch-and-bound, plans de coupe, relaxation lin\'eaire.
  \item[Chapitre 10 -- Programmation non lin\'eaire.]
    Conditions de Karush-Kuhn-Tucker, programmation convexe, m\'ethodes de
    gradient.
  \item[Chapitre 11 -- Simulation et th\'eorie des files d'attente.]
    Simulation \`a \'ev\'enements discrets, mod\`eles M/M/1, M/M/c, loi de Little.
\end{description}

\section*{Pr\'erequis}

Les pr\'erequis n\'ecessaires sont~:
\begin{itemize}
  \item \textbf{Alg\`ebre lin\'eaire}~: matrices, syst\`emes d'\'equations
        lin\'eaires, bases et rang.
  \item \textbf{Analyse}~: fonctions de plusieurs variables, d\'eriv\'ees
        partielles, convexit\'e.
  \item \textbf{Probabilit\'es}~: lois classiques, esp\'erance, processus de
        Poisson (pour le chapitre~11).
  \item \textbf{Algorithmique}~: notions de base en complexit\'e et en
        programmation Python.
\end{itemize}

\section*{Conventions et notations}

Tout au long de cet ouvrage, nous utiliserons les conventions suivantes~:

\begin{center}
\renewcommand{\arraystretch}{1.3}
\begin{tabular}{cl}
  \toprule
  \textbf{Notation} & \textbf{Signification} \\
  \midrule
  $\R$            & Ensemble des r\'eels \\
  $\R^n$          & Espace euclidien de dimension $n$ \\
  $\Z$            & Ensemble des entiers relatifs \\
  $\N$            & Ensemble des entiers naturels \\
  $\mathbf{x}$    & Vecteur colonne de $\R^n$ \\
  $A \in \R^{m \times n}$ & Matrice r\'eelle $m \times n$ \\
  $\mathbf{c}^\top \mathbf{x}$ & Produit scalaire \\
  $\norm{\mathbf{x}}$ & Norme euclidienne \\
  $\abs{S}$       & Cardinal de l'ensemble $S$ \\
  $\argmin$, $\argmax$ & Argument du minimum / maximum \\
  \bottomrule
\end{tabular}
\end{center}

Les th\'eor\`emes, d\'efinitions et propositions sont num\'erot\'es par chapitre.
Les exemples sont signal\'es par un encadr\'e et les exercices figurent en fin de
chapitre.

\begin{remarque}
Les encadr\'es \textbf{Remarque} fournissent des compl\'ements ou des mises en
garde. Les encadr\'es \textbf{Bonne pratique} offrent des conseils
m\'ethodologiques pour la mod\'elisation et l'impl\'ementation.
\end{remarque}

\section*{Logiciels utilis\'es}

Les exemples num\'eriques sont trait\'es en \textbf{Python~3} avec les
biblioth\`eques suivantes~:

\begin{itemize}
  \item \texttt{scipy.optimize.linprog} -- r\'esolution de programmes lin\'eaires~;
  \item \texttt{PuLP} -- mod\'elisation et r\'esolution de PL et PLNE~;
  \item \texttt{NetworkX} -- graphes, plus courts chemins, flots~;
  \item \texttt{SimPy} -- simulation \`a \'ev\'enements discrets~;
  \item \texttt{NumPy}, \texttt{Matplotlib} -- calcul num\'erique et
        visualisation.
\end{itemize}

\begin{codeblock}[title=\textbf{Installation des d\'ependances}]
\begin{minted}{python}
# pip install numpy scipy matplotlib pulp networkx simpy
import numpy as np
from scipy.optimize import linprog
import pulp
import networkx as nx
\end{minted}
\end{codeblock}

\section*{Comment utiliser ce cours}

Chaque chapitre suit une structure r\'ecurrente~:
\begin{enumerate}
  \item \textbf{Motivation} -- un probl\`eme concret qui introduit le besoin
        th\'eorique.
  \item \textbf{Th\'eorie} -- d\'efinitions, th\'eor\`emes et d\'emonstrations.
  \item \textbf{Algorithmes} -- description formelle et illustration sur un
        exemple.
  \item \textbf{Impl\'ementation} -- code Python comment\'e.
  \item \textbf{Exercices} -- probl\`emes de difficult\'e croissante
        ($\star$ \`a $\star\star\star$).
\end{enumerate}

\begin{bonnepratique}[title=\textbf{Apprentissage actif}]
La recherche op\'erationnelle s'apprend en \textit{pratiquant}. Il est fortement
recommand\'e de~:
\begin{itemize}
  \item Refaire les exemples \`a la main avant de consulter la solution.
  \item Impl\'ementer les algorithmes en Python.
  \item Chercher les exercices avant de regarder les corrig\'es.
  \item Explorer des variantes des probl\`emes propos\'es.
\end{itemize}
\end{bonnepratique}

\section*{Bref historique}

\begin{itemize}
  \item \textbf{1937}~: Kantorovich formule les premiers probl\`emes de
        programmation lin\'eaire pour la planification \'economique.
  \item \textbf{1939--1945}~: Les \'equipes de recherche op\'erationnelle
        britanniques et am\'ericaines optimisent les op\'erations militaires
        (convois, radars, bombardements).
  \item \textbf{1947}~: George Dantzig publie la m\'ethode du simplexe.
  \item \textbf{1951}~: Kuhn et Tucker \'enoncent les conditions d'optimalit\'e
        pour la programmation non lin\'eaire.
  \item \textbf{1956}~: Ford et Fulkerson proposent l'algorithme de flot
        maximal.
  \item \textbf{1958}~: Gomory introduit les coupes enti\`eres.
  \item \textbf{1960}~: Land et Doig formalisent le branch-and-bound.
  \item \textbf{1979}~: Khachiyan montre que la PL est r\'esoluble en temps
        polynomial (m\'ethode de l'ellipso\"ide).
  \item \textbf{1984}~: Karmarkar publie la m\'ethode des points int\'erieurs.
  \item \textbf{Ann\'ees 1990--2020}~: D\'eveloppement des solveurs commerciaux
        (CPLEX, Gurobi) et open source (GLPK, CBC), g\'en\'eralisation dans
        l'industrie.
\end{itemize}

\section*{R\'ef\'erences principales}

\begin{enumerate}
  \item \textsc{Hillier} F.S., \textsc{Lieberman} G.J.,
        \textit{Introduction to Operations Research}, McGraw-Hill,
        11\textsuperscript{e}~\'ed., 2021.
  \item \textsc{Bertsimas} D., \textsc{Tsitsiklis} J.N.,
        \textit{Introduction to Linear Optimization}, Athena Scientific, 1997.
  \item \textsc{Wolsey} L.A.,
        \textit{Integer Programming}, Wiley, 2\textsuperscript{e}~\'ed., 2020.
  \item \textsc{Chv\'atal} V.,
        \textit{Linear Programming}, W.H.~Freeman, 1983.
  \item \textsc{Minoux} M.,
        \textit{Programmation math\'ematique~: th\'eorie et algorithmes},
        Lavoisier, 2\textsuperscript{e}~\'ed., 2008.
  \item \textsc{Sakarovitch} M.,
        \textit{Optimisation combinatoire~: graphes et programmation lin\'eaire},
        Hermann, 1984.
\end{enumerate}

\section*{Remerciements}

Je remercie les coll\`egues et \'etudiants qui ont contribu\'e \`a l'am\'elioration
de ce cours par leurs remarques et suggestions. Toute erreur restante est de
ma seule responsabilit\'e.

\vfill
\begin{flushright}
\textit{Bonne lecture et bonne optimisation~!}
\end{flushright}
